\textbf{结论名称：} 三角形形状判定（用正弦/余弦定理边角互化）\\
\textbf{适用条件：} 在 \(\triangle ABC\) 中，已知边角混合关系式（如 \(a\cos A = b\cos B\)），需判定三角形形状。\\
\textbf{变量说明：} \(a,b,c\) 为内角 \(A,B,C\) 的对边，\(R\) 为外接圆半径。\\
\textbf{核心公式：}
利用正弦定理 \(a=2R\sin A,\, b=2R\sin B\) 边化角，常用判定链：
\[
\boxed{a\cos A = b\cos B \iff \sin A\cos A = \sin B\cos B \iff \sin 2A = \sin 2B \;\Longrightarrow\; A=B \ \text{或}\ A+B=\frac{\pi}{2}}
\]
\textbf{使用场景：} 解三角形时出现形如 \(a\cos A = b\cos B\)、\(a\sin A = b\sin B\) 等边角混合方程，判定等腰、直角或等腰直角三角形。\\
\textbf{简单例子：} 在 \(\triangle ABC\) 中，若 \(a\cos A = b\cos B\)，判断其形状。\\
解：由正弦定理得 \(\sin A\cos A = \sin B\cos B\)，即 \(\sin 2A = \sin 2B\)。\\
因 \(0<A,B<\pi\)，得 \(2A=2B\) 或 \(2A+2B=\pi\)，即 \(A=B\) 或 \(A+B=\frac{\pi}{2}\)。\\
故该三角形为等腰三角形或直角三角形。